\chapter{Capacité, décroissance et encombrement}\label{ch:rs:capacity-decay-and-crowding}

Du lundi 6 au jeudi 9 août 2007, le facteur valeur du marché boursier américain a perdu 2.28~\% et le facteur momentum 3.43~\%, soit dix et douze fois
leurs écarts-types quotidiens, tandis que le marché lui-même a progressé de 1.4~\%. Aucune nouvelle ne concernait les actions. Le vendredi, le facteur valeur a regagné 1.42~\%
et le momentum 1.15~\%. Khandani et Lo attribuent l'épisode au \omterm{def:rs:capacity-decay-and-crowding:crowding}{débouclage} de portefeuilles long--short actions que de nombreux fonds quantitatifs détenaient en commun, et au
retrait de leur capital par les teneurs de marché au moment même. La \omterm{def:rs:capacity-decay-and-crowding:capacity}{capacité} d'une stratégie n'est pas une propriété de son seul signal~: elle dépend de ses propres coûts, qui croissent
avec sa taille et de tous les autres acteurs qui négocient les mêmes titres. Ce chapitre trace des \omterm{def:rs:capacity-decay-and-crowding:capacity}{courbes de capacité} avec \texttt{firm.capacity}, étudie la décroissance de l'alpha en production,
et simule un \omterm{def:rs:capacity-decay-and-crowding:crowding}{débouclage}.

\section{Courbes de capacité}

\begin{definition}[Capacité d'une stratégie, courbe de capacité, taille maximisant le profit]\label{def:rs:capacity-decay-and-crowding:capacity}
Une \emph{courbe de capacité}\index{courbe de capacité} donne le rendement net, le ratio de Sharpe net et le profit en dollars d'une stratégie en fonction de la taille du capital qu'elle gère, chaque
taille étant négociée avec ses coûts. La \emph{taille maximisant le profit}\index{taille maximisant le profit} est la taille pour laquelle le profit en dollars culmine. La
\emph{capacité}\index{capacité} d'une stratégie est la taille au-delà de laquelle un critère énoncé n'est plus satisfait~: ici, la taille à laquelle son ratio de Sharpe net tombe à la moitié de son meilleur niveau.
\end{definition}

\begin{proposition}[Profit et taille en présence de coûts concaves]\label{prop:rs:capacity-decay-and-crowding:size}
Si le rendement brut $g$ d'une stratégie ne dépend pas de sa taille $A$, et que ses coûts, rapportés au capital, croissent comme $k\sqrt A$ (transactions proportionnelles au capital,
impact en racine carrée), son profit en dollars $A(g - k\sqrt A)$ est maximal en $A^* = (2g/3k)^2$, où son rendement net vaut $g/3$.
\end{proposition}

\begin{proof}
La dérivée $g - \frac32 k\sqrt A$ s'annule en $\sqrt{A^*} = 2g/3k$, et le rendement net y vaut $g - k\sqrt{A^*} = g/3$.
\end{proof}

La \omterm{def:rs:capacity-decay-and-crowding:capacity}{taille maximisant le profit} est donc plus grande que celle qui maximise le rendement du capital, laquelle est la plus petite~; Berk et Green ont bâti une théorie de la
gestion active fondée sur ce rendement d'échelle décroissant. La \Cref{fig:rs:capacity-decay-and-crowding:capacity} trace le portefeuille de cinquante titres du chapitre~27 à dix tailles, de
10~millions de dollars à 30~milliards de dollars, en deux versions. La règle fixe (coûts ignorés, prévision lissée sur cinquante jours) négocie de la même façon à chaque taille~: son ratio de Sharpe net
passe de 1.79 pour 10~millions de dollars à 0.34 pour 5~milliards de dollars et $-0.28$ pour 10~milliards de dollars, se réduit de moitié à 2.02~milliards de dollars, et son profit en dollars culmine à
3~milliards de dollars, 146~millions de dollars par an. Le portefeuille réoptimisé à chaque taille, avec les coûts pris en compte et un lissage d'un jour, culmine à un ratio de Sharpe de 2.43 pour 1~milliard de dollars
et reste à 1.36 pour 30~milliards de dollars, où son profit, 806~millions de dollars par an, croît encore~: réoptimiser en fonction de la taille est le moyen le plus efficace d'étendre
la \omterm{def:rs:capacity-decay-and-crowding:capacity}{capacité}, parce qu'elle change la quantité qu'un portefeuille négocie à mesure qu'il grossit.

\begin{omfigure}
\begin{tikzpicture}
\begin{semilogxaxis}[name=a,width=6.8cm,height=5.4cm,xlabel={\small taille du fonds (\$)},ylabel={\small ratio de Sharpe net},tick label style={font=\footnotesize},
  xmin=7e6,xmax=4e10,ymin=-2.1,ymax=2.8,grid=major,grid style={gray!25},legend pos=south west,legend style={font=\scriptsize},table/col sep=comma]
\addplot[black!60,thick,mark=o] table[x=size,y=sr_fixed]{figdata/research/28-capacity-decay-and-crowding/capacity.csv};
\addplot[omDef,very thick,mark=*] table[x=size,y=sr_opt]{figdata/research/28-capacity-decay-and-crowding/capacity.csv};
\legend{règle fixe,réoptimisé}
\end{semilogxaxis}
\begin{semilogxaxis}[at={(a.east)},anchor=west,xshift=1.3cm,width=6.8cm,height=5.4cm,xlabel={\small taille du fonds (\$)},ylabel={\small profit (millions de dollars par an)},
  tick label style={font=\footnotesize},xmin=7e6,xmax=4e10,ymin=-320,ymax=850,grid=major,grid style={gray!25},table/col sep=comma]
\addplot[black!60,thick,mark=o] table[x=size,y=profit_fixed]{figdata/research/28-capacity-decay-and-crowding/capacity.csv};
\addplot[omDef,very thick,mark=*] table[x=size,y=profit_opt]{figdata/research/28-capacity-decay-and-crowding/capacity.csv};
\end{semilogxaxis}
\end{tikzpicture}

\omcaption{\omterm{def:rs:capacity-decay-and-crowding:capacity}{Courbes de capacité} du portefeuille de cinquante titres du chapitre~27~: ratio de Sharpe net (à gauche) et profit en dollars (à droite) en fonction de la taille du fonds, pour la règle fixe et le portefeuille
réoptimisé à chaque taille (le profit inférieur à $-300$ millions de dollars est tracé à $-300$). Données~: \texttt{rs\_capacity.curve}.}\label{fig:rs:capacity-decay-and-crowding:capacity}
\end{omfigure}

\section{Décroissance de l'alpha en production}

La \omterm{def:rs:capacity-decay-and-crowding:capacity}{capacité} d'une stratégie change après sa mise en production, pour des raisons que le chapitre~13 a mesurées~: le signal décroît à mesure que d'autres le découvrent, et le déclin post-publication de
des facteurs publiés en est la trace publique. La \omterm{def:rs:capacity-decay-and-crowding:capacity}{capacité} décroît doublement, puisque le rendement brut baisse et que les coûts, désormais partagés avec les concurrents qui négocient
les mêmes titres, augmentent. La \omterm{def:rs:capacity-decay-and-crowding:capacity}{courbe de capacité} doit donc être retracée avec le \omterm{def:rs:anatomy-of-a-predictor:ic}{coefficient d'information} observé en production et avec les coûts mesurés sur ses propres exécutions
(chapitres~23 et~27), et une stratégie doit être dimensionnée en dessous de sa \omterm{def:rs:capacity-decay-and-crowding:capacity}{taille maximisant le profit} estimée, là où la courbe est plate~: le profit perdu en étant plus petite est
faible, la protection contre une courbe qui s'est déplacée est grande.

\section{Détecter l'encombrement}

\begin{definition}[Encombrement, position encombrée, comomentum, débouclage]\label{def:rs:capacity-decay-and-crowding:crowding}
L'\emph{encombrement}\index{encombrement} (crowding) désigne la présence de nombreux investisseurs sur les mêmes positions~; une \emph{position encombrée}\index{position encombrée} (crowded trade) est une position détenue par assez
d'entre eux pour que leur sortie conjointe fasse bouger les cours. Le \emph{comomentum}\index{comomentum} (Lou et Polk) est la corrélation moyenne des rendements anormaux des
titres détenus par une stratégie de momentum, mesure du nombre d'intervenants qui agissent sur eux. Un \emph{débouclage}\index{débouclage} (unwind) est la réduction rapide d'une position encombrée par certains
de ses détenteurs, ce qui fait bouger les cours auxquels les autres valorisent les leurs.
\end{definition}

Stein a nommé le problème~: pour une stratégie dont la valeur n'a aucun ancrage dans les fondamentaux, un arbitragiste ne peut savoir combien de ses pairs prennent la même position,
et l'effet de levier que chacun choisit en privé crée un risque de vente forcée pour tous. Pedersen a analysé pourquoi les traders se regroupent et pourquoi ils se précipitent vers la sortie. Trois indicateurs rendent l'\omterm{def:rs:capacity-decay-and-crowding:crowding}{encombrement}
visible (le \cref{lst:rs:capacity-decay-and-crowding:curve})~: la corrélation entre les rendements des titres propres à une stratégie, nette des facteurs du modèle de risque (\omterm{def:rs:capacity-decay-and-crowding:crowding}{comomentum}
pour un portefeuille momentum~: des valeurs élevées annonçaient des retournements)~; le recouvrement du portefeuille de la maison avec les estimations des portefeuilles des autres (à partir des déclarations de détention réglementaires, ou à partir
des rendements des produits des concurrents)~; et l'encours vendeur à découvert, comme les jours de couverture, côté vendeur (livre~1, chapitre~16). Aucun n'est concluant, et chacun est bruité~; ils
relèvent du tableau de bord des risques, à côté de la liquidité du portefeuille.

\section{Août 2007}

Le \omterm{def:rs:capacity-decay-and-crowding:crowding}{débouclage} du chapitre (le \cref{lst:rs:capacity-decay-and-crowding:unwind}) place cinq fonds de 2~milliards de dollars chacun dans des portefeuilles long--short de 500 actions, soit quatre fois leur
capital en exposition brute, une centaine de titres de chaque côté~; chaque action s'échange pour 200~millions de dollars par jour, avec une volatilité quotidienne de 2~\%. Chaque fonds classe les actions selon un mélange
d'un score commun et d'un score propre, avec une pondération qui fixe le degré de recouvrement de leurs portefeuilles. Un fonds vend une fraction de son portefeuille sur trois jours. Chaque jour, ses transactions
font bouger les cours selon un impact en racine carrée, dont 30~\% persiste tandis que le reste décroît avec une demi-vie d'un jour, et les quatre autres fonds valorisent leurs portefeuilles aux
cours ainsi déplacés.

\begin{center}
\small
\begin{tabular}{@{}lrrrr@{}}
\toprule
 & \multicolumn{4}{c}{recouvrement des portefeuilles (cosinus)}\\
\cmidrule(l){2-5}
part du portefeuille vendue & 0.04 & 0.18 & 0.43 & 1.00\\
\midrule
10~\% & $-0.03~\%$ & $-0.18~\%$ & $-0.42~\%$ & $-0.97~\%$\\
25~\% & $-0.05~\%$ & $-0.28~\%$ & $-0.66~\%$ & $-1.54~\%$\\
50~\% & $-0.08~\%$ & $-0.40~\%$ & $-0.93~\%$ & $-2.17~\%$\\
100~\% & $-0.11~\%$ & $-0.56~\%$ & $-1.32~\%$ & $-3.07~\%$\\
\bottomrule
\end{tabular}
\end{center}

Le tableau donne la perte maximale des autres fonds, en fraction de leur capital~; elle survient le troisième jour, le dernier de la vente, dans tous les cas. La perte croît
proportionnellement au recouvrement et avec la racine carrée de la part vendue, puisque l'impact est une racine carrée~: vendre un quart du portefeuille coûte aux autres moitié
autant que de tout vendre. Avec des portefeuilles identiques et la totalité du portefeuille vendue, les autres perdent 3.07~\% de leur capital sans négocier, et le vendeur 2.00~\% (il paie aussi
pour ses transactions). Puis l'\omterm{def:rs:transaction-costs-in-research:impact}{impact temporaire} se résorbe~: au quinzième jour, la perte des autres est de 1.30~\%, soit la part permanente (la \cref{fig:rs:capacity-decay-and-crowding:unwind}).
La reprise a lieu le vendredi d'août 2007, et ceux qui ont perdu le plus sont ceux qui ont vendu au plus bas, figeant la part temporaire. Un fonds qui sait que la
perte est temporaire et a le capital nécessaire pour tenir ne perd que la part permanente~; un fonds doté d'un stop-loss devient le vendeur suivant, et c'est ainsi que se propage un \omterm{def:rs:capacity-decay-and-crowding:crowding}{débouclage}.

\begin{omfigure}
\begin{tikzpicture}
\begin{axis}[width=11cm,height=5.8cm,xlabel={jours depuis le début de la vente},ylabel={perte des autres fonds (\% du capital)},tick label style={font=\small},
  label style={font=\small},xmin=0,xmax=15,ymin=-3.3,ymax=0.2,grid=major,grid style={gray!25},legend columns=4,
  legend style={font=\footnotesize,at={(0.5,-0.25)},anchor=north,draw=none,fill=none,/tikz/every even column/.append style={column sep=8pt}},
  table/col sep=comma]
\addplot[omDef!35,thick] table[x=day,y=f10]{figdata/research/28-capacity-decay-and-crowding/unwind_paths.csv};
\addplot[omDef!60,thick] table[x=day,y=f25]{figdata/research/28-capacity-decay-and-crowding/unwind_paths.csv};
\addplot[omDef!85,thick] table[x=day,y=f50]{figdata/research/28-capacity-decay-and-crowding/unwind_paths.csv};
\addplot[omThm,very thick] table[x=day,y=f100]{figdata/research/28-capacity-decay-and-crowding/unwind_paths.csv};
\draw[black!50,densely dotted] (axis cs:3,-3.3) -- (axis cs:3,0.2);
\legend{10~\% vendus,25~\% vendus,50~\% vendus,tout vendu}
\end{axis}
\end{tikzpicture}

\omcaption{La perte moyenne des quatre fonds qui n'ont pas vendu, aux livres identiques, lorsqu'un fonds vend une partie de son livre sur trois jours (la ligne
pointillée marque le dernier jour de vente). La part permanente de l'impact est ce qui subsiste. Données~: \texttt{rs\_capacity.unwind}.}\label{fig:rs:capacity-decay-and-crowding:unwind}
\end{omfigure}

\section{Tutoriel~: la semaine du 6 août 2007}\label{tut:rs:capacity-decay-and-crowding:tut}

\begin{tutorial}
\textbf{Objectif.} Tracer les \omterm{def:rs:capacity-decay-and-crowding:capacity}{courbes de capacité} du livre du chapitre~27, trouver la \omterm{def:rs:capacity-decay-and-crowding:capacity}{taille maximisant le profit} et la taille à laquelle le ratio de Sharpe net est divisé par deux, et simuler
le \omterm{def:rs:capacity-decay-and-crowding:crowding}{débouclage} sur une grille de recouvrements et de fractions vendues. \textbf{État final~:} le tableau, les \cref{fig:rs:capacity-decay-and-crowding:capacity,fig:rs:capacity-decay-and-crowding:unwind},
ainsi que les \omterm{def:rs:risk-models:families}{rendements des facteurs} d'août 2007 issus de \texttt{data/research/ff\_aug2007.csv}.
\begin{tutsteps}
\item \textbf{Fonctions auxiliaires de \omterm{def:rs:capacity-decay-and-crowding:capacity}{capacité}}~: la courbe, son sommet et la taille correspondant à une fraction du meilleur.
\omcode{code/firm/capacity/firm_capacity.py}{29}{46}{Courbe de capacité, taille maximisant le profit et taille de division par deux.}\label{lst:rs:capacity-decay-and-crowding:curve}
\item \textbf{Le \omterm{def:rs:capacity-decay-and-crowding:crowding}{débouclage}}~: transactions à l'ouverture, impact en racine carrée, parts permanente et décroissante, valorisation.
\omcode{code/firm/capacity/firm_capacity.py}{64}{89}{Fonds valorisant des livres qui se recouvrent pendant qu'un fonds vend.}\label{lst:rs:capacity-decay-and-crowding:unwind}
\item \textbf{Exécuter} \texttt{rs\_capacity.curve}, \texttt{summary}, \texttt{unwind(theta, fraction)}, \texttt{fig\_capacity.py} et \texttt{rs\_fetch\_aug2007.py} une seule fois.
\end{tutsteps}
\textbf{Pour aller plus loin.} Donner aux fonds qui ne vendent pas un stop-loss à 1~\% et les laisser vendre lorsqu'il est atteint~: la cascade~; rendre la part permanente dépendante de
la présence ou non des teneurs de marché (le retrait de capital décrit par Khandani et Lo).
\end{tutorial}

\section{Construction~: capacité et encombrement}\label{bld:rs:capacity-decay-and-crowding:capacity}

\begin{build}
\bfield{Rôle} Chaque stratégie est dimensionnée à partir de sa \omterm{def:rs:capacity-decay-and-crowding:capacity}{courbe de capacité} et surveillée pour l'\omterm{def:rs:capacity-decay-and-crowding:crowding}{encombrement}~; le risque de \omterm{def:rs:capacity-decay-and-crowding:crowding}{débouclage} par d'autres acteurs constitue un scénario dans les
tests de résistance de la société.
\bfield{Interface} \texttt{capacity\_curve(sizes, net\_returns, vols)}, \texttt{profit\_maximising(curve)}, \texttt{size\_at\_fraction(curve, frac, key)},
\texttt{avg\_pairwise\_corr(R)}, \texttt{comomentum(R\_long, R\_short)}, \texttt{overlap(w1, w2)},
\texttt{unwind(books, capitals, seller, fraction, days, adv, sigma, eta, permanent, half\_life, horizon)}.
\bfield{Règles} Les \omterm{def:rs:capacity-decay-and-crowding:capacity}{courbes de capacité} sont tracées à partir de \omterm{def:rs:vectorised-backtests:backtest}{backtests} relancés à chaque taille avec les coûts estimés de la société, et retracées avec les \omterm{def:rs:anatomy-of-a-predictor:ic}{coefficients d'information} en direct~; une
stratégie fonctionne en dessous de sa \omterm{def:rs:capacity-decay-and-crowding:capacity}{taille maximisant le profit}~; les surveillances de l'\omterm{def:rs:capacity-decay-and-crowding:crowding}{encombrement} sont revues avec la liquidité du livre.
\bfield{Tests d'acceptation} \texttt{code/firm/capacity/tests/}~: la \omterm{def:rs:capacity-decay-and-crowding:capacity}{taille maximisant le profit} et la taille de division par deux d'une stratégie aux coûts en racine carrée, comparées à leurs
formes closes~; corrélation moyenne par paires et \omterm{def:rs:capacity-decay-and-crowding:crowding}{comomentum} d'un facteur commun introduit~; recouvrement calculé à la main~; un \omterm{def:rs:capacity-decay-and-crowding:crowding}{débouclage} dans lequel le même livre perd et le livre opposé gagne
de même montant, l'\omterm{def:rs:transaction-costs-in-research:impact}{impact temporaire} qui revient, et l'\omterm{def:rs:transaction-costs-in-research:impact}{impact permanent} qui ne revient pas.
\bfield{Pour aller plus loin} Cascades avec stop-loss~; \omterm{def:rs:capacity-decay-and-crowding:crowding}{encombrement} estimé à partir des déclarations réglementaires de détention~; \omterm{def:rs:capacity-decay-and-crowding:capacity}{capacité} partagée entre les stratégies de la société.
\end{build}

\begin{omsources}
\item A.~E.~Khandani et A.~W.~Lo, ``What happened to the quants in August 2007? Evidence from factors and transactions data'', \emph{Journal of Financial Markets}
14(1), 2011.
\item D.~Lou et C.~Polk, ``Comomentum: inferring arbitrage activity from return correlations'', \emph{Review of Financial Studies} 35(7), 2022.
\item L.~H.~Pedersen, ``When everyone runs for the exit'', NBER Working Paper 15297, 2009.
\item J.~C.~Stein, ``Presidential address: sophisticated investors and market efficiency'', \emph{Journal of Finance} 64(4), 2009.
\item J.~B.~Berk et R.~C.~Green, ``Mutual fund flows and performance in rational markets'', \emph{Journal of Political Economy} 112(6), 2004.
\item K.~R.~French, Data Library, Fama/French 3 Factors and Momentum Factor (daily).
\end{omsources}

\section{Exercices}

\begin{exercise}[$\star$]\label{exo:rs:capacity-decay-and-crowding:1}
Une stratégie rapporte 10~\% par an brut et ses coûts atteignent 10~\% à 10~milliards de dollars, en croissant comme la racine carrée de la taille. Calculer sa \omterm{def:rs:capacity-decay-and-crowding:capacity}{taille maximisant le profit} et son rendement net
à cette taille.
\end{exercise}

\begin{exercise}[$\star$]\label{exo:rs:capacity-decay-and-crowding:2}
Pourquoi la \omterm{def:rs:capacity-decay-and-crowding:capacity}{taille maximisant le profit} est-elle plus grande que la taille maximisant le rendement du capital~?
\end{exercise}

\begin{exercise}[$\star$]\label{exo:rs:capacity-decay-and-crowding:3}
Dans le \omterm{def:rs:capacity-decay-and-crowding:crowding}{débouclage} du chapitre aux livres identiques, quelle est l'ampleur de la poussée de cours d'un jour sur un titre en position longue lorsque tout le livre est vendu sur trois jours~?
\end{exercise}

\begin{exercise}[$\star\star$]\label{exo:rs:capacity-decay-and-crowding:4}
Pourquoi la perte des autres fonds croît-elle avec la racine carrée de la fraction vendue, et proportionnellement au recouvrement~?
\end{exercise}

\begin{exercise}[$\star\star$]\label{exo:rs:capacity-decay-and-crowding:5}
Pendant la semaine du 6 août 2007, le marché a monté tandis que les facteurs valeur et momentum ont chuté de dix écarts-types. Qu'est-ce que cela révèle sur les acteurs qui échangeaient~?
\end{exercise}

\begin{exercise}[$\star\star$]\label{exo:rs:capacity-decay-and-crowding:6}
Pourquoi réoptimiser à chaque taille étend-il la \omterm{def:rs:capacity-decay-and-crowding:capacity}{capacité} bien davantage que toute modification du signal~?
\end{exercise}

\begin{exercise}[$\star\star\star$]\label{exo:rs:capacity-decay-and-crowding:7}
\emph{Programmation.} Avec \texttt{rs\_capacity.unwind}, comparer la perte des autres fonds au quinzième jour à leur perte maximale, et mettre ce rapport en relation avec la part permanente.
\end{exercise}

\begin{exercise}[$\star\star\star$]\label{exo:rs:capacity-decay-and-crowding:8}
\emph{Trouver la faille.} «~La \omterm{def:rs:capacity-decay-and-crowding:capacity}{capacité} de notre stratégie est de 5~milliards de dollars~: à cette taille, son \omterm{def:rs:vectorised-backtests:backtest}{backtest} affiche encore un ratio de Sharpe de 0.34 après coûts.~»
\end{exercise}

\section{Problème~: la semaine du 6 août 2007}

\begin{problem}[{Problème du week-end --- quelle taille, et avec qui}]\label{pb:rs:capacity-decay-and-crowding:1}
Le livre du chapitre~27 à de nombreuses tailles, et le \omterm{def:rs:capacity-decay-and-crowding:crowding}{débouclage} du chapitre.

\medskip\noindent\textbf{Partie I --- \omterm{def:rs:capacity-decay-and-crowding:capacity}{Capacité}.}
\begin{enumerate}
\item Énoncer et démontrer la \omterm{def:rs:capacity-decay-and-crowding:capacity}{taille maximisant le profit} sous des coûts en racine carrée.
\item Donner le ratio de Sharpe net de la règle à taille fixe, de 10~millions à 30~milliards de dollars.
\item Où son profit atteint-il son sommet, et où son ratio de Sharpe est-il divisé par deux~?
\item Que change la réoptimisation à chaque taille~?
\item Comment dimensionner la stratégie, et pourquoi en dessous de sa \omterm{def:rs:capacity-decay-and-crowding:capacity}{taille maximisant le profit}~?
\end{enumerate}

\medskip\noindent\textbf{Partie II --- Décroissance et \omterm{def:rs:capacity-decay-and-crowding:crowding}{encombrement}.}
\begin{enumerate}[resume]
\item Pourquoi la \omterm{def:rs:capacity-decay-and-crowding:capacity}{capacité} décroît-elle une fois la stratégie mise en production~?
\item Quelles sont les deux complications de Stein~?
\item Qu'est-ce que le \omterm{def:rs:capacity-decay-and-crowding:crowding}{comomentum}, et que prédisaient les valeurs élevées~?
\item Quelles surveillances placer sur le tableau de bord~?
\end{enumerate}

\medskip\noindent\textbf{Partie III --- Le \omterm{def:rs:capacity-decay-and-crowding:crowding}{débouclage}.}
\begin{enumerate}[resume]
\item Décrire les fonds simulés et le modèle d'impact.
\item Donner les pertes maximales des autres fonds sur la grille.
\item Quand le maximum est-il atteint, et que subsiste-t-il au quinzième jour~?
\item Qu'ont fait les facteurs valeur et momentum du 6 au 10 août 2007~?
\item Qu'ont identifié Khandani et Lo~?
\end{enumerate}

\medskip\noindent\textbf{Partie IV --- Le verdict.}
\begin{enumerate}[resume]
\item Énoncer le \emph{résultat nommé}~: la perte maximale des fonds qui ne vendent pas, en fonction de la fraction du portefeuille chevauchant qui est vendue, et
la taille à partir de laquelle le ratio de Sharpe net de la stratégie est divisé par deux.
\item Que doit faire un fonds au troisième jour d'un \omterm{def:rs:capacity-decay-and-crowding:crowding}{débouclage} qu'il n'a pas lancé~?
\item Comment les stop-loss modifieraient-ils le tableau~?
\item Comment l'ensemble des stratégies propres à la société génère-t-il lui-même de l'\omterm{def:rs:capacity-decay-and-crowding:crowding}{encombrement} (le netting du chapitre~27)~?
\item Quelles questions poser à l'auteur d'une nouvelle stratégie au sujet de sa \omterm{def:rs:capacity-decay-and-crowding:capacity}{capacité}~?
\item En une phrase, dire ce qui détermine la \omterm{def:rs:capacity-decay-and-crowding:capacity}{capacité} d'une stratégie.
\end{enumerate}
\end{problem}

\section{Questions d'entretien}

\begin{interviewq}[$\star$ \iqroles{chercheur, risque}]\label{iq:rs:capacity-decay-and-crowding:1}
Que s'est-il passé pour les fonds actions quantitatifs en août 2007~?
\end{interviewq}

\begin{interviewq}[$\star\star$ \iqroles{chercheur}]\label{iq:rs:capacity-decay-and-crowding:2}
Comment estimeriez-vous la \omterm{def:rs:capacity-decay-and-crowding:capacity}{capacité} d'une stratégie~?
\end{interviewq}

\begin{interviewq}[$\star\star$ \iqroles{chercheur, risque}]\label{iq:rs:capacity-decay-and-crowding:3}
Comment sauriez-vous si votre stratégie est encombrée~?
\end{interviewq}

\begin{interviewq}[$\star\star$ \iqroles{risque}]\label{iq:rs:capacity-decay-and-crowding:4}
Un autre fonds au portefeuille semblable au vôtre est en cours de liquidation. Que faites-vous~?
\end{interviewq}

\begin{interviewq}[$\star\star$ \iqroles{chercheur, trader}]\label{iq:rs:capacity-decay-and-crowding:5}
Pourquoi le profit en dollars d'une stratégie peut-il continuer d'augmenter alors que son ratio de Sharpe a commencé à baisser~?
\end{interviewq}

\begin{interviewq}[$\star\star\star$ \iqroles{chercheur}]\label{iq:rs:capacity-decay-and-crowding:6}
Modéliser la perte d'un fonds qui détient un portefeuille chevauchant celui d'un vendeur, avec un impact en racine carrée et une part permanente, puis en déduire comment cette perte varie avec la fraction vendue.
\end{interviewq}
